\chapter{Persamaan orde pertama} \label{fo:chapter}

%%%%%%%%%%%%%%%%%%%%%%%%%%%%%%%%%%%%%%%%%%%%%%%%%%%%%%%%%%%%%%%%%%%%%%%%%%%%%%

\section{Integral sebagai solusi}
\label{integralsols:section}

\sectionnotes{1 pertemuan (atau kurang)\EPref{, \S1.2 dalam \cite{EP}}\BDref{,
dibahas dalam \S1.2 dan \S2.1 di \cite{BD}}}

Suatu ODE orde pertama adalah persamaan berbentuk
\begin{equation*}
\frac{dy}{dx} = f(x,y) ,
\end{equation*}
atau cukup
\begin{equation*}
y' = f(x,y) .
\end{equation*}
Secara umum, tidak ada rumus atau prosedur sederhana yang dapat diikuti untuk
menemukan solusi.
Dalam beberapa pertemuan berikutnya, kita akan melihat kasus-kasus ketika solusi
tidak sulit diperoleh.
Pada bagian ini, kita meninjau kasus ketika $f$ hanya bergantung pada $x$,
yaitu persamaannya berbentuk
\begin{equation} \label{ias:inteq}
y' = f(x) .
\end{equation}
Kita dapat mengintegralkan (mencari antiturunan) kedua ruas persamaan terhadap $x$:
\begin{equation*}
\int y'(x) \,dx = \int f(x) \,dx + C ,
\end{equation*}
yakni,
\begin{equation*}
y(x) = \int f(x) \,dx + C .
\end{equation*}
Fungsi $y(x)$ ini sebenarnya adalah solusi umum.
Jadi, untuk menyelesaikan \eqref{ias:inteq},
kita mencari suatu antiturunan dari $f(x)$
lalu menambahkan konstanta sembarang untuk memperoleh solusi umum.

\begin{example}
Carilah solusi umum dari $y' = 3 x^2$.

Dari kalkulus dasar, kita tahu
bahwa solusi umum haruslah $y = x^3 + C$. Mari kita periksa dengan
menurunkannya:
$y' = 3x^2$. Kita memperoleh kembali \emph{persis} persamaan semula.
\end{example}

Kini saat yang tepat untuk membahas satu hal mengenai
notasi dan istilah dalam kalkulus. Buku-buku teks kalkulus
memperkeruh persoalan dengan membicarakan integral terutama sebagai sesuatu
yang disebut \myindex{integral tak tentu}. Integral tak tentu
sebenarnya adalah \emph{\myindex{antiturunan}}
(bahkan, seluruh keluarga antiturunan
dengan satu parameter). Sesungguhnya hanya ada satu integral, yaitu
integral tentu.
Satu-satunya alasan penggunaan notasi integral tak tentu adalah bahwa kita selalu dapat
menuliskan antiturunan sebagai integral tentu. Artinya, menurut teorema dasar
kalkulus, kita selalu dapat menuliskan
$\int f(x) \,dx + C$ sebagai
\begin{equation*}
\int_{x_0}^x f(t) \,dt + C .
\end{equation*}
Karena itulah digunakan istilah \emph{\myindex{mengintegralkan}}, padahal yang sebenarnya kita maksud mungkin
\emph{\myindex{mencari antiturunan}}.
Pengintegralan hanyalah salah satu cara untuk menghitung
antiturunan (dan cara ini selalu berhasil; lihat
contoh-contoh berikut). Pengintegralan didefinisikan sebagai luas di bawah grafik---kebetulan
cara ini juga menghitung antiturunan.
Demi konsistensi, kita akan tetap menggunakan
notasi integral tak tentu ketika menginginkan antiturunan,
dan Anda harus \emph{selalu} memandang integral tentu
sebagai cara untuk menuliskannya.

Biasanya, kita juga mempunyai kondisi awal seperti $y(x_0) = y_0$
untuk sepasang bilangan $x_0$ dan $y_0$ ($x_0$ sering kali 0, tetapi tidak selalu).
Dengan demikian, kita dapat menuliskan solusi secara elegan sebagai integral tentu.
Misalkan masalah kita adalah $y' = f(x)$, $y(x_0) = y_0$. Maka solusinya adalah
\begin{equation} \label{int:eqdef}
y(x) = \int_{x_0}^x f(t) \,dt + y_0 .
\end{equation}
Mari kita periksa!
Dengan teorema dasar kalkulus, kita memperoleh
$y' = f(x)$, dan ternyata $y$ memang merupakan suatu
solusi. Apakah solusi ini memenuhi kondisi awal? Ya,
$y(x_0) = \int_{x_0}^{x_0} f(t)\,dt + y_0 = y_0$. Ternyata benar!
